\scanpage{157}
\mspage{31}
% archive sheet number: 156

\[
pq=\mathrm{id}_B,\quad qj=i,
\]
ce qui prouve c$'$), puisque $i\in\Phi$.

\uline{Corollaire} 1: Soit $(\Phi,\Psi)$ un pré-couple de Quillen.
Alors on a
\[
\begin{aligned}
\widetilde{\Phi}=\Psi^*&=\text{ens. des objets de }\underline{\mathrm{Fl}}\,M\text{ qui sont}\\
&\phantom{=}\text{facteurs directs d'éléments de }\Phi,
\end{aligned}
\tag*{(4.3.)}
\]
et dualement
\[
\begin{aligned}
\overline{\Psi}=\Phi_*&=\text{ens. des objets de }\underline{\mathrm{Fl}}\,M\text{ qui sont}\\
&\phantom{=}\text{facteurs directs d'éléments de }\Psi.
\end{aligned}
\tag*{(4.4)}
\]

\uline{Corollaire 2} : Les conditions suivantes sur le
précouple de Quillen $(\Phi,\Psi)$ sont équivalentes

a) Le précouple est \uline{clos} i.e \deleted{c'est un couple de Quillen.} \inserted{$\Phi=\widetilde\Phi$, $\Psi=\overline\Psi$}

b) $\Phi=\Psi^*$, $\Psi=\Phi_*$

c) $\Phi$ et $\Psi$ sont stables par facteurs directs.

Soit maintenant
\[
W\subset\mathrm{Fl}\,M
\tag*{(4.5.)}
\]
un localiseur, satisfaisant les conditions habituelles
\[
\left\{\begin{aligned}
&\text{a) }W\text{ contient les iso}\\
&\text{b) si }x\overset f{\unsure{-|\longrightarrow}}y\overset g{\unsure{-|\longrightarrow}}z\text{ dans }M,\text{ si }\deleted{\unsure{gf}}\text{ deux parmi }f,g,gf\\
&\qquad\text{sont \unsure{des|dans} }W,\text{ alors \unsure{la|le} troisième aussi}
\end{aligned}\right.
\tag*{(4.6.)}
\]
($W$ modérément saturée) \deleted{Rappelons} Supposons
de plus \unsure{donnés|données} \deleted{deux} \inserted{quatre autres} parties
\[
C,TC,\quad F,TF
\tag*{(4.7)}
\]
de $\mathrm{Fl}\,M$, satisfaisant les conditions suivantes
\[
\left\{\begin{aligned}
&\deleted{TC\subset C,\quad TF\subset F}\\
&\text{a) }(C,TF)\text{ et }(\unsure{TC|T},F)\text{ sont des précouples de Quillen}\\
&\text{b) }TC\subset C\cap W,\quad TF\subset F\cap W\\
&\deleted{\unsure{\text{c) si }p\in F\cap W}\text{ \unsure{et si} }p=\unsure{ji|pi},\text{ avec}}\\
&\qquad\deleted{\unsure{i\in TC,\ p\in TF}}.
\end{aligned}\right.
\tag*{(4.8.)}
\]
